
% ============================================================
% Oppgave 2
% ============================================================

\hypertarget{oppgave2}{}
\section*{Oppgave 2}
\addcontentsline{toc}{section}{Oppgave 2}
Det analoge signalet er gitt som:
\[
x(t)=2\cos(2\pi100t)
\]
og samplingsfrekvensen er:
\[
f_s=800~\mathrm{Hz}
\]
Samplingsperioden blir:
\[
T_s
=
\frac{1}{f_s}
=
\frac{1}{800}
\]
Et sample tas ved:
\[
t=nT_s
\]
Dette settes inn i uttrykket for det analoge signalet:
\[
x[n]
=
2\cos(2\pi100nT_s)
\]
Siden:
\[
T_s=\frac{1}{800}
\]
får vi:
\[
x[n]
=
2\cos
\left(
2\pi\frac{100}{800}n
\right)
\]
som kan forenkles til:
\[
\boxed{
x[n]
=
2\cos\left(\frac{\pi}{4}n\right)
}
\]
Den digitale vinkelfrekvensen er dermed:
\[
\hat{\omega}
=
\frac{\pi}{4}
\]
Antall samples per periode kan finnes fra:
\[
N
=
\frac{f_s}{f}
=
\frac{800}{100}
=
8
\]
Den digitale periodelengden er derfor:
\[
\boxed{N=8}
\]
For en periode, $n=0,\ldots,7$, får vi:
\[
\begin{array}{c|c}
n & x[n] \\
\hline
0 & 2 \\
1 & \sqrt{2} \\
2 & 0 \\
3 & -\sqrt{2} \\
4 & -2 \\
5 & -\sqrt{2} \\
6 & 0 \\
7 & \sqrt{2}
\end{array}
\]
Sekvensen for en periode blir dermed:
\[
\boxed{
x[n]
=
\{2,\sqrt{2},0,-\sqrt{2},-2,-\sqrt{2},0,\sqrt{2}\}
}
\]